\documentclass{beamer}
\usetheme{Madrid}
\usecolortheme{default}
\usepackage{graphicx, hyperref, xcolor}
\usepackage[warn]{mathtext}
\usepackage[T2A]{fontenc}
\usepackage[utf8]{inputenc}
\usepackage[english,russian]{babel}
\usepackage{indentfirst}
\usepackage{subcaption}
\usepackage{float}
\usepackage{amsmath}
\usepackage{mathrsfs}
\usepackage{amsfonts}
\usepackage{amssymb}
\usepackage{booktabs}
\usepackage{tikz}
\usetikzlibrary{arrows.meta, positioning}
\usepackage{graphicx}

\setbeamertemplate{navigation symbols}{}
\setbeamertemplate{footline}[frame number]

\title{Оптимизация ансамблей градиентного бустинга через адаптивное обучение агрегирующих функций}
\author{Якушевич Антон Сергеевич}
\institute{ВМК МГУ \\[1em] \textit{Научный руководитель}: д.ф.-м.н. Сенько Олег Валентинович}
\date{2025}


\begin{document}

\frame{\titlepage}

\begin{frame}{Актуальность}
    \begin{itemize}
        \item Градиентный бустинг — мощный метод в машинном обучении
        \item Проблема: выбор агрегирующей функции
        \item Цель: снизить количество деревьев в ансамбле, сохранив при этом качество модели
    \end{itemize}
\end{frame}

% --------------------
\begin{frame}{Постановка задачи}
    Пусть задана последовательность функций $\{F_N\}_{N=1}^{+\infty}$, где $F_N: \mathbb{R}^N \to \mathbb{R}$. На $(N+1)$-ом шаге обучения вместо стандартной суммы $\sum_{i=1}^{N} T_i(x)$ будем использовать обобщённую функцию ансамбля $F_N\left(T_1(x), T_2(x), \dots, T_N(x)\right)$, где $\{T_i\}_{i=1}^N$ -- базовые модели (решающие деревья), обученные на предыдущих шагах.

    \vspace{0.5cm}

    \textbf{Дано:}

    Выборка $X^\ell = \{(x_i, y_i)\}_{i=1}^\ell$, где $y_i \in \mathbb{R}$.
    
    \vspace{0.3cm}
    
    \textbf{Найти:}
    
    Последовательность агрегирующих функций $F_N$.
    
    \vspace{0.3cm}
    
    \textbf{Критерий:}
    
    Улучшение качества при фиксированном $N$ относительно функции потерь $\mathscr{L}$ ($RMSE$, $MAE$ и $R^2$), по сравнению с классическим градиентным бустингом, использующим сумму в качестве агрегирующей функции.
\end{frame}


% --------------------
\begin{frame}{Параметризация агрегирующей функции}
    Разложение в ряд Тейлора функции $F_N(T_1(x), T_2(x), \dots, T_N(x))$ будет выглядеть следующим образом:
    \[
        \sum_{k=0}^{+\infty} \Bigg[ \overbrace{\sum_{k_1=0}\sum_{k_2=0}\dots\sum_{k_N=0}}^{k_1 + k_2 + \dots + k_N = k} \frac{1}{k_1!k_2!\dots k_N!}\frac{\partial^k F_N(a_1, a_2, \dots, a_N)}{\partial T_1^{k_1} \partial T_2^{k_2} \dots \partial T_N^{k_N}} \prod_{i=1}^N (T_i(x)-a_i)^{k_i}\Bigg]
    \]

    В качестве агрегирующей функции будем использовать усечённый ряд Тейлора:
    \[
        F_N(x) = \sum_{k=1}^{M} \overbrace{\sum_{k_1=0}\sum_{k_2=0}\dots\sum_{k_N=0}}^{k_1 + k_2 + \dots + k_N = k} b_{k_1, k_2, \dots, k_N} T_1^{k_1}(x) T_2^{k_2}(x) \dots T_N^{k_N}(x)
    \]

\end{frame}

% --------------------
\begin{frame}{Стратегия обучения нового алгоритма}
    \textbf{1. Нахождение точного корня:}
    \[
        F_{N+1}(x_i) = y_i \quad \Rightarrow \quad \sum_{k=0}^M C_k(x_i) \cdot (T_{N+1}(x_i))^k = y_i
    \]

    \textbf{2. Градиент:}
    \[
        T_{N+1}(x_i) = -\alpha \cdot C_1(x_i) \cdot \frac{\partial\mathscr{L}(y_i, F)}{\partial F} \Big|_{F = F_N(x_i)}
    \]

    \begin{itemize}
        \item \footnotesize Learning rate: $\alpha = \frac{\delta}{N^\theta}$
        \item $C_k(x)$ — функции, однозначно задаваемые $T_1(x), \dots, T_N(x)$
    \end{itemize}
\end{frame}

% --------------------
\begin{frame}{Оптимизация коэффициентов $b$}
Для фиксированного датасета $\mathscr{D}$ вводится оператор $A_\mathscr{D} (b): \mathbb{R}^P \to \mathbb{R}$, который запускает кросс-валидацию: на каждом фолде обучает наш бустинг с заданными коэффициентами $b$ и возвращает среднюю ошибку по всем фолдам.

Частные производные:
\[
\frac{\partial A_\mathscr{D}}{\partial b_i}\Bigg|_{b=b_0}  \approx \frac{A_\mathscr{D}(b_0 + \varepsilon \cdot e_i) - A_\mathscr{D}(b_0)}{\varepsilon}
\]

Градиентный спуск с моментом:
\[
\begin{cases}
    v^{(k+1)} = \mu v^{(k)} - \eta \nabla_b A_\mathscr{D}(b^{(k)}) \\
    b^{(k+1)} = b^{(k)} + v^{(k+1)}
\end{cases}
\]

\end{frame}

% --------------------
\begin{frame}{Мета-модель}
\begin{center}
\scalebox{0.6}{
    \begin{tikzpicture}[
        node distance=1.5cm and 1cm,
        layer/.style={rectangle, draw=black!50, fill=black!10, minimum width=2.2cm, minimum height=0.8cm, rounded corners=2pt, align=center},
        input/.style={layer, fill=blue!10},
        output/.style={layer, fill=green!10},
        operator/.style={rectangle, draw=red!50, fill=red!10, minimum width=2.5cm, minimum height=1cm, rounded corners=2pt, align=center},
        arrow/.style={-Stealth, thick}
    ]
    
        % --- Первый ряд ---
        \node[input] (input) {Вход \\ $\mathscr{D}$};
        \node[layer, right=of input] (proj) {Input \\ projection};
        \node[layer, right=of proj] (encoder) {Encoder \\ TransformerBlocks};
        \node[layer, right=of encoder] (pma) {Pooling};
        \node[layer, right=of pma] (mlp) {MLP \\ + Residuals};
    
        % --- Второй ряд ---
        \node[output, below=1.8cm of mlp] (output) {Выход};
        \node[operator, left=of output] (operator) {$A_{\mathscr{D}}(b)$};
        \node[left=of operator] (loss) {(1) $\log(1 + \mathscr{L}) + \| b \|_2$};
        \node[right=of output] (sec_out) {$b$ (2)}
    
        % --- Стрелки первого ряда ---
        \draw[arrow] (input) -- (proj);
        \draw[arrow] (proj) -- (encoder);
        \draw[arrow] (encoder) -- (pma);
        \draw[arrow] (pma) -- (mlp);
    
        % --- Переход вниз и второй ряд ---
        \draw[arrow] (mlp.south) -- ++(0,-0.6) -| (output.north);
        \draw[arrow] (output) -- (operator);
        \draw[arrow] (output) -- (sec_out);
        \draw[arrow] (operator) -- (loss);
    
        % --- Skip-connection сверху ---
        \path (output) -- (operator) coordinate[midway] (mid1);
        \path (operator) -- (loss) coordinate[midway] (mid2);
        
        \draw[arrow]
            (mid1)
                |- ([yshift=0.7cm]operator.north)
                node[above] {$\| b \|_2$}
                -| (mid2);
    
    \end{tikzpicture}
}
\end{center}

Оптимизационная задача:
$$
\min_{\phi} \Big[ \log \Big(1 + A_{\mathscr{D}}(\Psi_{\phi}(\mathscr{D})) \Big) + \| \Psi_{\phi}(\mathscr{D}) \|_2 \Big]
$$
где $\Psi_{\phi}$ — нейронная сеть с параметрами $\phi$, формирующая вектор коэффициентов $b = \Psi_{\phi}(\mathscr{D})$, используемых при построении ансамбля

\end{frame}

% --------------------
\begin{frame}{Модели для сравнения}
    \begin{itemize}
        \item \textbf{GB}: Стандартный градиентный бустинг (сумма деревьев).
        \item \textbf{FGBReg} (Functional Gradient Boosting Regressor): Базовый вариант предлагаемой модели с следующей формулой агрегирующей функции:
            \[
                F_N(x) = \sum_{k=1}^{M} \sum_{k_1 + \dots + k_N = k} \frac{1}{k} \cdot T_1^{k_1}(x) \cdot T_2^{k_2}(x) \cdots T_N^{k_N}(x)
            \]
        \item \textbf{FGBReg + GD}: Оптимизация коэффициентов градиентным спуском.
        \item \textbf{FGBReg + NN}: Оптимизация с помощью нейронной мета-модели.
    \end{itemize}
\end{frame}

% --------------------
\begin{frame}{Результаты}

\begin{table}[H]
\centering
\begin{tabular}{|c|c|c|c|c|}
    \hline
     Набор данных & GB              & FGBReg   & FGBReg + GD    & FGBReg + NN     \\
    \hline
    task1         & 93.69           & 95.72   & \textbf{86.78}  & 99.43           \\
    task6         & 123.8           & 119.7   & \textbf{107.4}  & 132.8           \\
    task7         & 8.424           & 8.380   & \textbf{8.366}  & 8.378           \\
    \hline
    task10        & 0.177           & 0.178   & 0.233           & \textbf{0.161}  \\
    task16        & 41.03           & 33.54   & \textbf{32.05}  & 35.99           \\
    task18        & \textbf{68.85}  & 77.60   & 74.45           & 81.50           \\
    \hline
    task41        & 0.401           & 0.364   & 0.356           & \textbf{0.345}  \\
    task901       & \textbf{309.0}  & 361.7   & 367.1           & 655.7           \\
    task970       & \textbf{40688}  & 46121   & 45923           & 51427           \\
    \hline
\end{tabular}
\caption{$RMSE$}
\label{tab:rmse_results}
\end{table}
    
\end{frame}

\begin{frame}{Результаты}

\begin{table}[H]
\centering
\begin{tabular}{|c|c|c|c|c|}
    \hline
     Набор данных & GB              & FGBReg           & FGBReg + GD    & FGBReg + NN     \\
    \hline
    task1         & 67.09           & 72.21           & \textbf{62.40}  & 79.40           \\
    task6         & 78.33           & 89.09           & \textbf{76.50}  & 107.2           \\
    task7         & 6.857           & 6.788           & \textbf{6.786}  & 6.816           \\
    \hline
    task10        & 0.165           & 0.166           & 0.222           & \textbf{0.151}  \\
    task16        & 32.83           & 25.95           & \textbf{24.84}  & 28.70           \\
    task18        & \textbf{51.55}  & 57.69           & 53.00           & 64.39           \\
    \hline
    task41        & 0.211           & \textbf{0.194}  & 0.197           & 0.202           \\
    task901       & \textbf{189.1}  & 244.5           & 249.6           & 424.6           \\
    task970       & \textbf{28994}  & 32799           & 32810           & 34680           \\
    \hline
\end{tabular}
\caption{$MAE$}
\label{tab:mae_results}
\end{table}
    
\end{frame}

\begin{frame}{Результаты}

\begin{table}[H]
\centering
\begin{tabular}{|c|c|c|c|c|}
    \hline
     Набор данных & GB              & FGBReg           & FGBReg + GD     & FGBReg + NN     \\
    \hline
    task1         & 0.704           & 0.711            & \textbf{0.755}  & 0.695           \\
    task6         & 0.654           & 0.723            & \textbf{0.762}  & 0.661           \\
    task7         & 0.319           & 0.335            & 0.336           & \textbf{0.340}  \\
    \hline
    task10        & -0.863          & 0.207            & -0.108          & \textbf{0.303}  \\
    task16        & -11.28          & \textbf{-10.42}  & -11.61          & -13.33          \\
    task18        & \textbf{0.681}  & 0.642            & 0.603           & 0.617           \\
    \hline
    task41        & -0.503          & -0.199           & -0.133          & \textbf{-0.057} \\
    task901       & \textbf{0.919}  & 0.894            & 0.891           & 0.716           \\
    task970       & \textbf{0.728}  & 0.695            & 0.694           & 0.634           \\
    \hline
\end{tabular}
\caption{$R^2$}
\label{tab:r2_results}
\end{table}
    
\end{frame}

% --------------------
\begin{frame}{Выводы}
\begin{itemize}
    \item Заменили сумму на гибкую агрегирующую функцию.
    \item Научились оптимизировать её коэффициенты.
    \item При тех же 20 деревьях оптимизация (особенно GD) заметно снижает $RMSE$/$MAE$ и повышает $R^2$
    \item Мета-сеть быстро генерирует хорошие коэффициенты для новых задач
\end{itemize}
\end{frame}

\begin{frame}[allowframebreaks]{Список литературы}
\tiny
\nocite{*}
\bibliographystyle{unsrtnat}
\bibliography{references}
\end{frame}


\end{document}
